% SPDX-License-Identifier: Apache-2.0
\section{A Channel Between Two Clocks}
\label{sec:x-cdc}

Two clocks that are not the same clock have no fixed relation. An
edge of one falls wherever it falls with respect to the other, and a
register written by the first and read by the second is read at an
arbitrary moment in its life, including the moment it is changing.
What comes out then is neither the old value nor the new one, and the
design has read something that never existed.

\code{ChanCdc} in \code{//lib/parts} carries words across that gap.
It is one unit with two processes, one per clock, a memory written by
the writing side and read by the reading side, and a pointer on each
side saying how far it has got. What crosses is only the pointers.

\subsection{Why the pointers are Gray coded}

A binary pointer changes several bits at once. Seven to eight is
\code{0111} to \code{1000}, four bits, and a sample caught in the
middle of that can read any of sixteen values, most of which the
pointer never held. A reader that believes one of them concludes the
queue is a different length than it is, and the error is unbounded.

A Gray code changes one bit per step, so a sample caught mid-change
reads either the value before or the value after and nothing else.
Both are safe, and they are safe in a particular direction: the
reader concludes the queue is emptier than it is and the writer that
it is fuller, so each is wrong in the direction that refuses work
rather than the one that corrupts it. A pointer sampled stale costs a
cycle; a pointer sampled impossible costs the data.

The pointers are one bit wider than the address so that full and
empty can be told apart, which the same number of bits cannot do:
equal pointers would mean both. Empty is the two equal. Full is the
two equal in every bit but the top two, which is what one lap apart
looks like in Gray.

\lstinputlisting[rangebeginprefix=//\ begin\{,rangebeginsuffix=\},%
rangeendprefix=//\ end\{,rangeendsuffix=\},includerangemarker=false,%
linerange=state-state,caption={\code{lib/parts/src/cdc.rs}: the
state the two processes share. Each side owns its own pointer and
keeps the other's through two flops of its own clock.},%
label={lst:x-cdc}]%
{lib/parts/src/cdc.rs}

\subsection{What the check covers, and what it cannot}

The example runs the writing side on a clock of period four and the
reading side on one of period six with a phase of two, so the edges
drift past each other and every relative position happens somewhere
in the run. The build checks the netlist against that run under
\code{nvc} and under Verilator, at the ports and at every register:
both Gray pointers and all four synchroniser flops are compared, so a
pointer that advanced at the wrong moment, or a full test off by the
wrap bit, is a failure rather than a difference nobody looks at.

What it does not cover is the reason the two flops are there.
Metastability is not modelled and cannot be: both sides are
deterministic, and the sampling clock reads a register that has
settled. The agreement demonstrated here is about the pointer
protocol. The second flop is an argument about silicon, and no
simulation in either language says anything about it. A green test
beside a synchroniser invites the opposite conclusion, which is why
this paragraph is here.

Writing this example found a fault in the checking rather than in the
design: a channel's \code{ready} and \code{valid} were compared
against the trace at ticks which were not edges of the port's own
clock, where the trace holds nothing for them. Under the default
clock every tick the testbench visits is an edge, so it had never
shown. That is issue 384, and it is the reason no example before this
one had a channel on a slower clock.

\begin{figure*}[htbp]
\centering
\input{cdc_timing.tex}
\caption{Words crossing from \code{clk\_w} to \code{clk\_r}.
\code{wgray} steps one bit at a time, \code{0,1,3,2,6,7,5,4},
appears at \code{wgray\_r2} two reading edges later, and the word
follows. \code{inp\_ready} is the wire whose checking was wrong
until issue 384.}
\label{fig:x-cdc}
\end{figure*}

\subsection{What it does not replace}

\code{lib/board/chan\_cdc.v} is the hand-written Verilog this is
shaped after, and the flagship holds five instances of it. Whether
those become this is a separate decision and a separate change: that
file is proven on hardware and this one is proven in simulation, and
those are not the same claim.
